% Generated draft: requires mathematical and source review.
\documentclass[11pt]{article}
\usepackage[margin=25mm]{geometry}
\usepackage{fontspec}
\setmainfont{Latin Modern Roman}
\usepackage{amsmath,amssymb,mathtools,unicode-math}
\setmathfont{Latin Modern Math}
\usepackage{enumitem,needspace,xurl,hyperref}
\setlength{\parindent}{0pt}
\setlength{\parskip}{5pt}
\setlength{\emergencystretch}{3em}
\newcommand{\N}{\mathbb{N}}
\newcommand{\Z}{\mathbb{Z}}
\newcommand{\Q}{\mathbb{Q}}
\newcommand{\R}{\mathbb{R}}
\newcommand{\C}{\mathbb{C}}
\newcommand{\F}{\mathbb{F}}
\newcommand{\A}{\mathbb{A}}
\newcommand{\PP}{\mathbb{P}}
\newcommand{\E}{\mathbb{E}}
\newcommand{\eps}{\varepsilon}
\newcommand{\dd}{\,\mathrm{d}}
\begin{document}
\section*{UnsolvedMath ID: 20002903}
\textbf{Source problem code:} catalogue entry
\subsection{Definitions and mathematical statement}
Work in the ring of symmetric functions over $\Q(q)$. The source asks to describe expansions of LLT polynomials in $k$-Schur functions, as a possible route to Schur expansions, and notes that certain LLT polynomials are already $k$-Schur functions. Neither the LLT normalization, the graded $k$-Schur convention, nor the rule choosing $k$ is specified in the quoted problem.

In the supplied new-variant convention, for an ordered tuple of skew diagrams $\boldsymbol\nu=(\nu^{(0)},\ldots,\nu^{(\ell-1)})$ use semistandard fillings and adjusted cell contents $\widetilde c(u)=\ell c(u)+i$ for $u\in\nu^{(i)}$, where $c$ is column minus row. An attacking ordered pair satisfies $0<\widetilde c(v)-\widetilde c(u)<\ell$; it contributes an inversion when $T(u)>T(v)$. Write
\[
G_{\boldsymbol\nu}(X;q)=\sum_Tq^{\operatorname{inv}(T)}X^{\operatorname{wt}(T)}.
\]
Let $\omega$ be the standard involution $\omega s_\lambda=s_{\lambda'}$. In this convention the aid says the Haiman--Haglund expansion formulation uses $\omega G_{\boldsymbol\nu}$, not the raw $G_{\boldsymbol\nu}$, in the graded basis $s^{(k)}_\lambda(X;q)$, with the $k$-Schur grading parameter identified with $q$. An expansion
\[
\omega G_{\boldsymbol\nu}=\sum_{\substack{\lambda\vdash|\boldsymbol\nu|\\\lambda_1\leq k}}a^{(k)}_{\lambda,\boldsymbol\nu}(q)s^{(k)}_\lambda(X;q)
\]
requires membership in the corresponding $k$-Schur span; the record does not supply general membership hypotheses, an admissible $k$ rule, or a complete general coefficient formula.

For $n$ ordered single cells of equal ordinary content, all pairs attack and the aid gives
\[
G_{\boldsymbol\nu}(X;q)=\sum_{\mu\vdash n}\frac{[n]_q!}{\prod_j[\mu_j]_q!}\,m_\mu(X),
\qquad [a]_q!=\prod_{j=1}^a(1+q+\cdots+q^{j-1}),
\]
with $[0]_q!=1$, and $\omega G_{\boldsymbol\nu}=s^{(1)}_{(1^n)}(X;q)$. In degree two this reads
\[
G=s_{(2)}+q s_{(1,1)},\qquad
\omega G=s_{(1,1)}+q s_{(2)}=s^{(1)}_{(1,1)}(X;q).
\]
The raw $G$ is not in the one-dimensional degree-two graded $1$-Schur span over $\Q(q)$, so omitting $\omega$ already fails at span membership. This statement is over the rational-function field, not a claim about every numerical specialization of $q$.

Separately, if the attack graph of a tuple $(\lambda^{(i)}/\mu^{(i)})_{i=1}^\ell$ is empty, the aid gives $G=\prod_i s_{\lambda^{(i)}/\mu^{(i)}}$. For $k$ larger than the total degree $d$, the stable $k$-Schur expansion is its Schur expansion. Explicitly its coefficient at $s^{(k)}_\lambda=s_\lambda$ is
\[
\sum_{\nu^{(1)},\ldots,\nu^{(\ell)}}
\left(\prod_{i=1}^\ell c^{\lambda^{(i)}}_{\mu^{(i)},\nu^{(i)}}\right)c^\lambda_{\nu^{(1)},\ldots,\nu^{(\ell)}},
\]
where the $c$ are ordinary and iterated Littlewood--Richardson coefficients. This is a nonnegative explicit family, not a general LLT-to-$k$-Schur formula.
\subsection{Short English statement}
The LLT-to-$k$-Schur question depends on normalization and on which graded span is intended. The supplied two-cell example detects a missing $\omega$ involution before positivity can even be asked. It also supplies explicit equal-content single-cell and attack-free families, with stable expansions in the latter case. A general expansion rule is not provided.
\subsection{Sources}
\begin{itemize}
\item Local registry: permanent catalogue entry, exact source code catalogue entry; \texttt{aim-special-functions-notes.json}, \texttt{notes[7]}.
\item AIM workshop \emph{Generalized Kostka polynomials}, Problem (8), \url{https://aimath.org/WWN/kostka/prob2.pdf}. No proposer, year, original remarks, or original literature reference is supplied.
\item Supplied research attempt 1, registry updated 14 August 2026: partial result, low novelty confidence, machine-generated and not peer reviewed. It refers to a Haiman--Haglund expansion formulation without a title, date, or full authoritative citation, and supplies the explicit families and degree-two normalization obstruction.
\item The verification field contains only ``The exact AIM record is:'' and no fuller normalization check.
\end{itemize}
\end{document}
